% Part: normal-modal-logic
% Chapter: filtrations
% Section: finite

\documentclass[../../../include/open-logic-section]{subfiles}

\begin{document}

\olfileid{nml}{fil}{fin}

\olsection{వడపోతల పరిమితత్వం}

\tetoken{ఉపసూత్రాల}{!!{formula}s} పరంగా సంవృతమైన
ఏ సమితి~$\Gamma$కైనా వడపోతలను నిర్వచించాం.
ఆ నిర్వచనం ఒక్కటే వడపోతలు పరిమితమని హామీ ఇవ్వదు.
నిజానికి, $\Gamma$ అనంతమైతే (ఉదాహరణకు, అన్ని
\tetoken{సూత్రాల}{!!{formula}s} సమితి అయితే),
వడపోత కూడా అనంతం కావచ్చు. అయితే $\Gamma$
పరిమితమైతే (ఉదాహరణకు, ఇచ్చిన
\tetoken{సూత్రం}{!!{formula}}~$!A$ యొక్క
\tetoken{ఉపసూత్రాల}{!!{formula}s} సమితి అయితే),
$\Gamma$ ద్వారా చేసే ఏ వడపోతైనా పరిమితమే.

\begin{prop}\ollabel{prop:filt-are-finite}
  $\Gamma$ పరిమితమైతే, $\Gamma$ ద్వారా నమూనా
  $\mModel{M}$కు చేసే ఏ వడపోత $\mModel{M^*}$ అయినా
  పరిమితమే.
\end{prop}

\begin{proof}
  $W^*$ పరిమాణం, $\equiv$ తుల్యతా సంబంధం కింద
  ఏర్పడే వేర్వేరు వర్గాలు~$[w]$ ఎన్ని ఉన్నాయో
  అంతే. ఒక వర్గంలోని ఏ రెండు లోకాలు $u$, $v$ అయినా,
  అంటే $u \equiv v$ అయ్యే ఏ $u$, $v$ అయినా,
  $\Gamma$లోని అన్ని \tetoken{సూత్రాల}{!!{formula}s}~$!A$
  విషయంలో ఏకీభవిస్తాయి: $!A \in \Gamma$ అయితే
  $!A$ రెండు లోకాలు $u$, $v$ వద్దా సత్యం, లేదా రెండింటి
  వద్దా అసత్యం. అందువల్ల ప్రతి వర్గం~$[w]$
  $\Gamma$ యొక్క ఒక ఉపసమితికి అనురూపం; అది
  $[w]$లోని లోకాల వద్ద $!A$ సత్యమయ్యే అన్ని
  $!A \in \Gamma$ల సమితి. వేర్వేరు వర్గాలు
  $[u]$, $[v]$ ఒకే $\Gamma$ ఉపసమితికి అనురూపం
  కావు. ఎందుకంటే $u$ వద్ద సత్యమయ్యే
  \tetoken{సూత్రాల}{!!{formula}s} సమితి, $v$ వద్ద
  సత్యమయ్యే \tetoken{సూత్రాల}{!!{formula}s} సమితి
  ఒకటే అయితే, $u$, $v$లు~$\Gamma$లోని అన్ని
  సూత్రాలపై ఏకీభవిస్తాయి; అంటే $u \equiv v$.
  అప్పుడు $[u] = [v]$. కాబట్టి $W^*$ నుంచి
  $\Pow{\Gamma}$కు
  \tetoken{ఒక ఒకటి-ఒకటి}{!!a{injective}} ప్రమేయం
  ఉంది; అందువల్ల $\card{W^*} \le
  \card{\Pow{\Gamma}}$. $\Gamma$లో $n$ వాక్యాలు
  ఉంటే, $W^*$ కార్డినాలిటీ~$2^n$ను మించదు.
\end{proof}

\end{document}
